\documentclass[11pt]{article}
\usepackage[a4paper,margin=25mm]{geometry}
\usepackage{amsmath,amssymb,amsthm}
\usepackage[hidelinks]{hyperref}
\newtheorem*{proposition}{Counterexample}
\title{A self-dual monotone identity is not antisymmetric}
\author{Conjecture 00000008845 \quad ziangni-sys}
\date{}
\begin{document}
\maketitle
The source classifies self-dual monotone operators by densely defined antisymmetric extensions and antisymmetric bilinear forms. The identity on the real Hilbert line is a counterexample. We explicitly verify self-duality in the convex-Lagrangian sense, rather than relying on the undefined label; we also record its actual adjoint and inverse-graph identities. Antisymmetry here is the skew symmetry of bilinear forms, as in the source's classification, not order-theoretic antisymmetry.

\begin{proposition}
On $H=\mathbb R$, the bounded linear operator $T(x)=x$ is maximal monotone and is generated by a proper continuous convex self-dual Lagrangian. Its graph has no antisymmetric extension. Moreover, $T$ cannot extend any densely defined antisymmetric relation.
\end{proposition}
\begin{proof}
Consider the everywhere finite quadratic Lagrangian
\[
 L(x,p)=\frac{x^2+p^2}{2}.
\]
It is continuous, nonnegative, and convex. More explicitly, for $a,b\geq0$, $a+b=1$, and $z,w\in\mathbb R^2$, its convexity gap is
\[
 aL(z)+bL(w)-L(az+bw)
 =\frac{ab}{2}\big((z_1-w_1)^2+(z_2-w_2)^2\big)\geq0.
\]
Its full Fenchel conjugate is computed over all $(u,v)\in\mathbb R^2$:
\[
 L^*(q,r)=\sup_{u,v}\{qu+rv-L(u,v)\}
 =\frac{q^2+r^2}{2}.
\]
Indeed the difference from this upper bound is
$-\tfrac12((u-q)^2+(v-r)^2)$, and equality is attained at $(q,r)$. Consequently $L^*(p,x)=L(x,p)$: this is genuine self-duality under the exchange of primal and dual variables. The standard equality graph generated by $L$ is
\[
 \{(x,p):L(x,p)=xp\}
 =\{(x,p):(x-p)^2=0\}
 =\{(x,x):x\in\mathbb R\}=\operatorname{gra}T.
\]
This computation includes all dual functionals: every continuous linear functional $f$ on $\mathbb R^2$ has the form $f(u,v)=f(1,0)u+f(0,1)v$.

Monotonicity is $(x-y)(Tx-Ty)=(x-y)^2\geq0$. To check maximality directly, let a monotone relation contain the identity graph and an additional point $(x,p)$. Test monotonicity against $(a,a)$ with $a=(x+p)/2$. It yields
\[
 0\leq(x-a)(p-a)=-\frac{(x-p)^2}{4},
\]
so $p=x$. Thus no larger monotone graph exists. The identity is also defined on all of $H$, satisfies the actual Hilbert adjoint identity $T^*=T$, and its inverse relation equals its graph.

For any antisymmetric relation $G$, skew symmetry requires
\[
 \langle p,y\rangle=-\langle x,q\rangle
 \quad\text{whenever }(x,p),(y,q)\in G.
\]
An extension of $\operatorname{gra}T$ contains $(1,1)$; using that point twice would give $1=-1$. Conversely, if an antisymmetric relation is contained in $\operatorname{gra}T$, each of its points is $(x,x)$ and the same condition gives $x^2=-x^2$, hence $x=0$. Its domain is contained in $\{0\}$ and is therefore not dense in $\mathbb R$. This rules out both directions of the terse extension wording.
\end{proof}

The associated form of the identity is $B(x,y)=xy$, which is symmetric and nonzero, not antisymmetric. Thus self-duality and monotonicity cannot imply the asserted skew classification. No failure of the usual theory of self-dual convex Lagrangians is asserted.

\paragraph{Formal verification.}
The Lean project uses the actual real Hilbert space, continuous linear identity and adjoint. It proves joint convexity and lower semicontinuity of $L$, and computes the Fenchel supremum for \emph{every} continuous linear functional on the product space, including attainment. It identifies the generated graph, proves its maximal monotonicity and inverse-graph identity, and excludes both antisymmetric extensions and dense-domain antisymmetric restrictions. The pinned Lean 4.19.0 project prints axiom audits of these results.
\end{document}
